%% kapa: exam=1 round=29 year=2018 period=1 form=A subject=civil_law subject_ko=민법 seq=5 number=5
%% source: Q-Net 감정평가사 기출문제 (cst00309) artlSeq=5202808 file=extracted/1008/5202808/2018년 제29회 감정평가사 1차시험 기출문제/2018년 제29회 감정평가사 1차시험 기출문제/1교시 A형.pdf
%% conversion: kordoc 4.24.51 -> tools/gen1.py (auto); review=proofread-sonnet; flags=none
\begin{question}{5}
법인의 불법행위책임에 관한 설명으로 옳지 않은 것은? (다툼이 있으면 판례에 따름)
\begin{choices}
  \item 대표자는 그 명칭이나 직위는 문제되지 않으며, 대표자로 등기되지 않은 자도 이에 포함될 수 있다.
  \item 대표자의 행위가 직무에 관한 것이 아님을 피해자가 안 경우, 법인은 책임을 지지 않는다.
  \item 외형상 대표자의 직무행위로 인정되어도 그것이 대표자 개인의 사리를 도모하기 위한 것이면, 직무에 관한 행위에 해당하지 않는다.
  \item 법인의 책임이 성립하는 경우 특별한 사정이 없는 한, 사원이 그 사항의 총회의 결에 찬성했다는 사실만으로 법인과 연대책임을 부담하지는 않는다. % [원문 그대로: 총회의 결에]
  \item 법인책임이 대표자의 고의적인 불법행위로 인한 경우에도 피해자에게 과실이 있다면, 법원은 이를 참작하여야 한다.
\end{choices}
\end{question}
\answer{3}
